\cotefile{11}
\scanpage{43}
\mspage{-3.7-}
% archive sheet number: 42
(identifying $\handwritten{\mathbb{A}}_R$ as $\handwritten{\mathbb{A}}^{n\unsure{\nu|v}-n}$), so that $(\handwritten{\mathbb{A}}^{n\nu-n},\handwritten{\phi}_R)$ is a geometric
quotient of $\handwritten{\unsure{\mathcal{U}|\mathscr{U}}}_R$ by $\mathrm{PGL}(n)$. Moreover, recall that
\[
\handwritten{\sigma}_R:\handwritten{\mathbb{A}}^{n\nu-n}\handwritten{\longrightarrow}\handwritten{\mathcal{U}}_R
\]
is a global section for $\handwritten{\phi}_R$. The following is obvious:

\uline{Corollary} \uline{3.2}: Suppose $\handwritten{\mathcal{U}\subset\mathcal{U}}_R$ is any invariant open subset. Then
$\handwritten{\mathcal{U}\cong}\mathrm{PGL}(n)\mathbin{\mathrm{x}}\handwritten{\sigma}_R^{-1}(\handwritten{\mathcal{U}})$, the isomorphism being compatible with the
obvious operations of $\mathrm{PGL}(n)$.

To illustrate these results, which may be somewhat obscured by the
cumbersome machinery of $R$-partitions, look at the case $\nu=1$, $S_1=\handwritten{\{}0,1,\ldots,n\handwritten{\}}$.
The proposition is then the classical result that $n+2$ points in
``general'' position in $\handwritten{\mathbb{P}}_n$ can be normalized to the canonical points
$(1,0,\ldots,0)$, $(0,1,\ldots,0),\ldots,(0,0,\ldots,1)$, and $(1,1,\ldots,1)$ by a unique
projective transformation. To give a second example, suppose $n=2$,
$\nu=2$,
\[
S_1=\handwritten{\{}0,1\handwritten{\}},\quad S_2=\handwritten{\{}1,2\handwritten{\}}.
\]
The proposition asserts that, given 5 points $P_0,P_1,P_2,P_3$ and $P_4$ in
the plane, suppose $P_0,P_1$ and $P_2$ are not collinear, that $P_3$ is not
on the line $\overline{P_0P_2}$ or the line $\overline{P_1P_2}$ and that $P_4$ is not on the line
$\overline{P_0P_1}$ or $\overline{P_0P_2}$: then we can uniquely normalize $P_0,P_1,P_2$ to $(1,0,0)$,
$(0,L,0)$ and $(0,0,1)$ respectively, and normalize $P_3$ to $(1,1,a)$, $P_4$
to $(b,1,1)$ by a suitable projective transformation. (cf. \handwritten{Figure 3},
where $Q_3=(1,1,0)$, $Q_4=(0,1,1)$ ).

